\scanpage{23}
à puissances divisées est un \inserted{pseudo-}\uline{anneau} à puissances divisées, muni d'une structure de $A$-algèbre, satisfaisant
\[
1.6.\qquad\pi^n(\lambda x)^{\unsure{n}}=\lambda^n\pi^n(x)\qquad\lambda\in A,\quad x\in\unsure{D|J}.
\]
\margin{N.B. C'est tjs vérifié si $\Lambda=\mathbb Z$}

Si $J$ est un idéal de $A$, c'est une $A$-algèbre; on \deleted{\unsure{dit que la structure à puissances divisées}\ \unsure{compatible}\ \illegible{2}\ \unsure{scalaires}} \inserted{appelle structure à puissances divisées sur \uline{l'idéal}} $J$ une structure de puissances divisées sur la \deleted{ps-anneau} \inserted{$A$-algèbre} sous-jacente, i.e. une structure de p.div. sur le pseudo-anneau sous-jacent, satisfaisant 6.

\uline{\S2} Soit $A$ un anneau, et $M$ un $A$-module. Considérons \inserted{la $A$-algèbre}
\[
\Gamma^*(M)=\coprod_{n\ge0}\Gamma^n(M),
\]
on a des applications
\[
\pi^n_{(1)}:x\rightsquigarrow x^{(n)}\qquad\text{de }M\text{ dans }\Gamma^nM
\]
On va les prolonger en des applications
\[
\pi^n:\Gamma^+(M)\longrightarrow\Gamma^+(M)\qquad x\rightsquigarrow x^{(n)}
\]
induisant
\[
\pi^n_{(p)}:\Gamma^p(M)\longrightarrow\Gamma^{np}(M)\qquad\deleted{\unsure{x\rightsquigarrow x^{(n)}}}
\]
en notant que ceci revient à se donner une ``application polynomiale homogène'' de degré $p$
\[
x\rightsquigarrow\pi^n_{(p)}(x^{(p)}):\underset{\sim}{M}\longrightarrow\underset{\sim}{\Gamma^{np}(M)}
\]
\inserted{``$(x^{(p)})^{(n)}$}
et on prendra
\[
(x^{(p)})^{(n)}=x^{(pn)}\frac{(pn)!}{(p!)^n n!}
\]
\mspage{22}
